
%% bare_conf.tex
%% V1.4b
%% 2015/08/26
%% by Michael Shell
%% See:
%% http://www.michaelshell.org/
%% for current contact information.
%%
%% This is a skeleton file demonstrating the use of IEEEtran.cls
%% (requires IEEEtran.cls version 1.8b or later) with an IEEE
%% conference paper.
%%
%% Support sites:
%% http://www.michaelshell.org/tex/ieeetran/
%% http://www.ctan.org/pkg/ieeetran
%% and
%% http://www.ieee.org/

%%*************************************************************************
%% Legal Notice:
%% This code is offered as-is without any warranty either expressed or
%% implied; without even the implied warranty of MERCHANTABILITY or
%% FITNESS FOR A PARTICULAR PURPOSE! 
%% User assumes all risk.
%% In no event shall the IEEE or any contributor to this code be liable for
%% any damages or losses, including, but not limited to, incidental,
%% consequential, or any other damages, resulting from the use or misuse
%% of any information contained here.
%%
%% All comments are the opinions of their respective authors and are not
%% necessarily endorsed by the IEEE.
%%
%% This work is distributed under the LaTeX Project Public License (LPPL)
%% ( http://www.latex-project.org/ ) version 1.3, and may be freely used,
%% distributed and modified. A copy of the LPPL, version 1.3, is included
%% in the base LaTeX documentation of all distributions of LaTeX released
%% 2003/12/01 or later.
%% Retain all contribution notices and credits.
%% ** Modified files should be clearly indicated as such, including  **
%% ** renaming them and changing author support contact information. **
%%*************************************************************************


% *** Authors should verify (and, if needed, correct) their LaTeX system  ***
% *** with the testflow diagnostic prior to trusting their LaTeX platform ***
% *** with production work. The IEEE's font choices and paper sizes can   ***
% *** trigger bugs that do not appear when using other class files.       ***                          ***
% The testflow support page is at:
% http://www.michaelshell.org/tex/testflow/



\documentclass[conference]{IEEEtran}
% Some Computer Society conferences also require the compsoc mode option,
% but others use the standard conference format.
%
% If IEEEtran.cls has not been installed into the LaTeX system files,
% manually specify the path to it like:
% \documentclass[conference]{../sty/IEEEtran}





% Some very useful LaTeX packages include:
% (uncomment the ones you want to load)


% *** MISC UTILITY PACKAGES ***
%
%\usepackage{ifpdf}
% Heiko Oberdiek's ifpdf.sty is very useful if you need conditional
% compilation based on whether the output is pdf or dvi.
% usage:
% \ifpdf
%   % pdf code
% \else
%   % dvi code
% \fi
% The latest version of ifpdf.sty can be obtained from:
% http://www.ctan.org/pkg/ifpdf
% Also, note that IEEEtran.cls V1.7 and later provides a builtin
% \ifCLASSINFOpdf conditional that works the same way.
% When switching from latex to pdflatex and vice-versa, the compiler may
% have to be run twice to clear warning/error messages.






\ifCLASSINFOpdf

\else

\fi


% correct bad hyphenation here
\hyphenation{op-tical net-works semi-conduc-tor}


\begin{document}
%
% paper title
% Titles are generally capitalized except for words such as a, an, and, as,
% at, but, by, for, in, nor, of, on, or, the, to and up, which are usually
% not capitalized unless they are the first or last word of the title.
% Linebreaks \\ can be used within to get better formatting as desired.
% Do not put math or special symbols in the title.
\title{Kişiselleştirilmiş Haber Öneri Sistemi için Özellik Vektörü Çıkarımı ve Analizi}


% author names and affiliations
% use a multiple column layout for up to three different
% affiliations
\author{\IEEEauthorblockN{Yunus Karatepe}
\IEEEauthorblockA{Bilgisayar Mühendisliği Bölümü \\
İstanbul Teknik Üniversitesi\\
İstanbul, Türkiye \\
karatepe22@itu.edu.tr}
}

% conference papers do not typically use \thanks and this command
% is locked out in conference mode. If really needed, such as for
% the acknowledgment of grants, issue a \IEEEoverridecommandlockouts
% after \documentclass

% for over three affiliations, or if they all won't fit within the width
% of the page, use this alternative format:
% 
%\author{\IEEEauthorblockN{Michael Shell\IEEEauthorrefmark{1},
%Homer Simpson\IEEEauthorrefmark{2},
%James Kirk\IEEEauthorrefmark{3}, 
%Montgomery Scott\IEEEauthorrefmark{3} and
%Eldon Tyrell\IEEEauthorrefmark{4}}
%\IEEEauthorblockA{\IEEEauthorrefmark{1}School of Electrical and Computer Engineering\\
%Georgia Institute of Technology,
%Atlanta, Georgia 30332--0250\\ Email: see http://www.michaelshell.org/contact.html}
%\IEEEauthorblockA{\IEEEauthorrefmark{2}Twentieth Century Fox, Springfield, USA\\
%Email: homer@thesimpsons.com}
%\IEEEauthorblockA{\IEEEauthorrefmark{3}Starfleet Academy, San Francisco, California 96678-2391\\
%Telephone: (800) 555--1212, Fax: (888) 555--1212}
%\IEEEauthorblockA{\IEEEauthorrefmark{4}Tyrell Inc., 123 Replicant Street, Los Angeles, California 90210--4321}}




% use for special paper notices
%\IEEEspecialpapernotice{(Invited Paper)}




% make the title area
\maketitle

% As a general rule, do not put math, special symbols or citations
% in the abstract
\begin{abstract}
\end{abstract}

% IEEE formatında anahtar kelimeler (isteğe bağlı ama önerilir)
% \begin{IEEEkeywords}
% Tavsiye Sistemleri, Veri Madenciliği, Özellik Mühendisliği, Haber Tavsiyesi, RecSys 2024, Gradient Boosting
% \end{IEEEkeywords}

% --- BÖLÜM 1: GİRİŞ VE PROBLEM TANIMI ---
\section{Giriş ve Problem Tanımı}
\label{sec:problem}


Dijital haber platformları, kullanıcıları platformda tutmak ve reklam gelirlerini maksimize etmek için kişiselleştirilmiş içerik sunumuna büyük önem vermektedir. Bu bağlamda, tavsiye sistemleri, her kullanıcıya ilgi alanlarına en uygun makaleleri sunarak kilit bir rol oynamaktadır.

Bu projenin temel problemi, kullanıcıya bir haber platformunda gösterilen makaleler arasından hangisine tıklayacağını yüksek doğrulukla tahmin eden bir sıralama modeli geliştirmektir. Bu, RecSys 2024 Mücadelesi'nin de temelini oluşturan bir ``tıklama oranı tahmini" (Click-Through Rate - CTR) problemidir.

Haber tavsiye sistemleri, e-ticaret platformlarından farklı olarak, içeriğin çok hızlı eskimesi (yüksek devir hızı) ve kullanıcı ilgilerinin anlık değişebilmesi gibi dinamik zorluklar barındırır. Bu dinamik ortamda, modelin başarısı büyük ölçüde veriden çıkarılan özelliklerin kalitesine ve amaca uygunluğuna bağlıdır. Bu proje, RecSys 2024 veri kümesinin sunduğu zengin bağlamsal, davranışsal ve içerik verilerinden yola çıkarak gelişmiş özellik mühendisliği (feature engineering) tekniklerine odaklanacaktır. Çalışmanın temel amacı sadece yüksek doğruluklu bir tıklama tahmini modeli geliştirmek değil, aynı zamanda hangi özelliklerin (örn. kullanıcının geçmiş okuma alışkanlıkları, makalenin popülerliği, anlık oturum davranışı) model başarısına ne kadar etki ettiğini nicel olarak analiz etmektir.


% --- BÖLÜM 2: KULLANILACAK VERİ KÜMESİ ---
\section{Kullanılacak Veri Kümesi}
\label{sec:dataset}

Projede, RecSys 2024 Mücadelesi için Danimarka haber sitesi Ekstra Bladet tarafından sağlanan EB-NeRD (Ekstra Bladet News Recommendation Dataset) veri kümesi kullanılacaktır. Bu veri kümesi, 6 haftalık bir periyottan toplanan gerçek kullanıcı davranışlarını içermektedir. Veri, `.parquet` formatında dört ana dosyadan oluşmaktadır:

\begin{enumerate}
    \item \textbf{behaviors.parquet:} Kullanıcıların hangi makaleleri gördüğünü (`Inview Article IDs') ve bunlardan hangilerine tıkladığını (`Clicked Article IDs') içeren ana etkileşim dosyasıdır. Kullanıcının cihaz tipi, o anki oturumdaki (session) okuma süresi gibi bağlamsal bilgileri de içerir.
    \item \textbf{history.parquet:} Kullanıcıların `behaviors' dosyasındaki etkileşimlerden önceki 21 günlük tıklama geçmişlerini (makale ID'leri, okuma süreleri, kaydırma yüzdeleri) içerir.
    \item \textbf{articles.parquet:} Veri kümesindeki tüm makalelerin meta verilerini içerir. Başlık, alt başlık, kategori, yayınlanma zamanı ve en önemlisi makale popülerliği (`Total Pageviews' vb.) gibi zengin içerik bilgileri sunar.
    \item \textbf{artifacts.parquet:} Makalelerin metin içeriklerinden (BERT, RoBERTa vb. ile) üretilmiş gömme vektör temsillerini (embeddings) içerir.
\end{enumerate}

Bu zengin veri yapısı, sadece işbirlikçi filtreleme (collaborative filtering) değil, aynı zamanda içerik tabanlı ve hibrit modellerin geliştirilmesine de olanak tanımaktadır.

% --- YENİ BÖLÜM: İLGİLİ ÇALIŞMALAR ---
\section{İlgili Çalışmalar}
\label{sec:relatedwork}

Haber tavsiye sistemleri literatürü, e-ticaretin aksine, içeriğin çok kısa ömürlü olması ve kullanıcı ilgi alanlarının yüksek oranda dinamik olması gibi zorluklara odaklanmıştır. Bu proje, RecSys 2024 mücadelesi için yayınlanan ve bu zorlukları içeren EB-NeRD veri kümesini \cite{kruse2024recsys} temel almaktadır.

Literatürdeki modern yaklaşımlar genellikle iki ana kategoriye ayrılmaktadır:

\textbf{1) Derin Öğrenme Tabanlı Kullanıcı Modelleme:} Bu yaklaşımlar, kullanıcıların okuma geçmişinden ve makale içeriklerinden zengin temsiller (embeddings) öğrenmeyi amaçlar. Örneğin, NPA \cite{wu2019npa}, kullanıcıların haber başlıklarındaki ve içeriklerindeki bilgilere farklı ağırlıklar vermesini sağlayan kişiselleştirilmiş dikkat (attention) mekanizmaları önermiştir. Benzer şekilde, LSTUR \cite{an2019lstur}, kullanıcıların hem uzun vadeli genel tercihlerini (tüm okuma geçmişinden) hem de kısa vadeli anlık ilgilerini (mevcut oturumdan) modelleyen çift bileşenli bir mimari sunmuştur. Bu modeller, özellikleri ``uçtan-uca" (end-to-end) öğrenir.

\textbf{2) Sıralı (Sequential) Modeller:} Özellikle kullanıcıların oturum (session) içi davranışlarına odaklanan bu modeller, tıklama sırasının önemini vurgular. GRU4Rec \cite{hidasi2015gru4rec}, RNN (Tekrarlayan Sinir Ağları) kullanarak oturumdaki bir sonraki tıklamayı tahmin eden temel çalışmalardan biridir. SASRec \cite{kang2018sasrec} ise bu alanda Transformer mimarisini kullanarak daha başarılı sonuçlar elde etmiştir.

Bu projede önerilen yaklaşım, literatürdeki bu çalışmalardan farklı bir noktaya odaklanmaktadır. Belirli bir model mimarisine odaklanmaktan ziyade, projenin ana katkısı \texttt{history}, \texttt{behaviors} ve \texttt{articles} dosyalarındaki zengin veriyi birleştirerek kapsamlı özellikler türetmek olacaktır. Çalışma, bu özelliklerin (örn. kullanıcı-makale benzerlik skorları, makale popülerliği, tazelik, bağlamsal oturum verileri) model başarısı üzerindeki etkisini incelemeyi amaçlamaktadır. Türetilen özellik setinin gücü, farklı modelleme yaklaşımlarından bağımsız olarak değerlendirilebilecek ve hangi veri kaynaklarının tıklama tahmini için daha değerli olduğuna dair analizler sunacaktır.


% --- YENİ BÖLÜM: ÖNERİLEN YAKLAŞIM ---
\section{Önerilen Yaklaşım}
\label{sec:approach}

Bu proje, literatürdeki uçtan-uca derin öğrenme modellerini ya da ardışık modelleri geliştirmek veya model eğitim aşamasında geliştirme sunmak yerine özellik vektörü çıkarımı aşamasında bir gelişim vaat etmektedir. Yapılacak çalışma temel olarak iki aşamalı gibi düşünülebilir: (1) Kapsamlı Özellik Mühendisliği ve (2) Model Eğitimi ve Özellik Analizi.

\subsection{Kapsamlı Özellik Mühendisliği}
Projenin temelini oluşturan bu aşamada, dört ana kategoride özellikler (feature) türetilecektir:

\textbf{1. İçerik ve Popülerlik Özellikleri:} \texttt{articles.parquet} dosyasından elde edilecek bu özellikler, makalenin ``tıklanabilirliğini" ölçmeyi hedefler.
\begin{itemize}
    \item \textbf{Popülerlik:} `Total Pageviews', `Total Inviews' ve `Tıklama Oranı (CTR)'.
    \item \textbf{Tazelik (Freshness):} Makalenin yayınlanma zamanı ile gösterim anı arasındaki saniye/saat farkı.
    \item \textbf{İçerik:} `Category String', `Premium' durumu, `Sentiment Score' gibi meta veriler.
\end{itemize}

\textbf{2. Kullanıcı Geçmişi Özellikleri:} \texttt{history.parquet} dosyasından yola çıkarak her kullanıcı için bir profil oluşturulacaktır.
\begin{itemize}
    \item \textbf{Kullanıcı İlgi Profili:} Kullanıcının geçmişte en çok okuduğu kategoriler ve konular (NER/Topics).
    \item \textbf{Kullanıcı Davranış Profili:} Kullanıcının ortalama `Read times' ve `Scroll Percentages` değerleri.
\end{itemize}

\textbf{3. Hibrit Kullanıcı-Makale Etkileşim Özellikleri:} Bu, projedeki en kritik özellik grubudur ve kullanıcı ile o anki aday makale arasındaki ``uyumu" ölçer.
\begin{itemize}
    \item \textbf{İçerik Benzerlik Skoru:} Kullanıcının okuma geçmişindeki makalelerin (`artifacts.parquet` vektörlerinin ağırlıklı ortalaması) oluşturduğu ``Kullanıcı Profili Vektörü" ile o an gösterimde olan ``Aday Makale Vektörü" arasındaki kosinüs benzerliği skoru.
    \item \textbf{Kategori Eşleşmesi:} Aday makalenin kategorisinin, kullanıcının geçmişte en çok okuduğu 3 kategori içinde olup olmadığı (binary özellik).
\end{itemize}

\textbf{4. Bağlamsal (Anlık) Özellikler:} \texttt{behaviors.parquet} dosyasındaki anlık durumdan alınacaktır.
\begin{itemize}
    \item `Device Type` (Cihaz Tipi), `Time` (günün saati, haftanın günü), `SSO Status`.
    \item \textbf{Oturum İçi Davranış:} Kullanıcının mevcut `Session ID` içindeki `Readtime` ve `Scroll Percentage` gibi anlık bağlılık (engagement) sinyalleri.
\end{itemize}

\subsection{Model Eğitimi ve Özellik Analizi}
Özellik çıkarımının ardından bu veri üzerinde, tıklama tahmini görevini gerçekleştirecek bir makine öğrenmesi modeli eğitilecektir. Bu aşamada, belirli bir model mimarisinden ziyade, projenin asıl odak noktası olan özelliklerin gücünü en iyi şekilde yansıtabilecek bir yaklaşım benimsenecektir. Farklı temel modeller (baseline) üzerinde denemeler yapılarak, özelliklerin katkısı ölçülecektir.

Model eğitiminden sonra özellik analizi yapılacaktır. Seçilen modelin sunduğu yorumlanabilirlik teknikleri veya modelden bağımsız (model-agnostic) analiz yöntemleri aracılığıyla, türetilen hangi özelliklerin tıklama tahminine pozitif veya negatif yönde ne kadar katkı sağladığı incelenecektir.










% --- BÖLÜM 5: DEĞERLENDİRME YÖNTEMİ ---
\section{Değerlendirme Yöntemi}
\label{sec:evaluation}

Modelin performansı, RecSys 2024 Mücadelesi'nde de kullanılan standart endüstri metrikleri ile ölçülecektir. Temel amaç bir sıralama problemi olduğu için, metrikler de sıralama kalitesine odaklanacaktır:

\begin{enumerate}
    \item \textbf{AUC (Area Under the Curve):} Modelin pozitif (tıklanan) ve negatif (tıklanmayan) örnekleri ne kadar iyi ayırt ettiğini ölçer. Genel sıralama kalitesi için standart bir metriktir.
    \item \textbf{MRR (Mean Reciprocal Rank):} Sadece ilk doğru tahminin pozisyonunu önemser. Listenin en başında doğru makaleyi bulmanın kritik olduğu durumlar için önemlidir.
    \item \textbf{nDCG@k (normalized Discounted Cumulative Gain):} Listenin ilk 'k' pozisyonundaki tüm doğru tahminleri, pozisyonlarına göre ağırlıklandırarak değerlendirir. Gerçek dünya sistemlerinde en çok tercih edilen metriklerden biridir.
\end{enumerate}

% --- BÖLÜM 6: ZAMAN PLANLAMASI VE B PLANI ---
\section{Zaman Planlaması ve B Planı}
\label{sec:plan}

Proje, belirlenen adımlara ve takvime göre yürütülecektir. Öngörülen zaman planlaması Tablo \ref{tab:zaman_planlamasi}'nde gösterilmiştir.

% Zaman Planlaması Tablosu
\begin{table}[h!]
  \caption{Öngörülen Proje Zaman Planlaması (8 Hafta)}
  \label{tab:zaman_planlamasi}
  \centering
  % Sütun formatı güncellendi (l: sol, p{...}: belirli genişlikte paragraf)
  \begin{tabular}{@{}l p{6cm}@{}} 
    \toprule
    \textbf{Hafta} & \textbf{Görev} \\
    \midrule
    Hafta 1 & Veri analizi, temizliği ve dosyaların birleştirilmeye hazırlanması. \\
    \addlinespace % Satırlar arası boşluk
    Hafta 2 & Temel model (baseline) kurulumu (örn. popülerlik özellikleri ile) ve ilk temel skorun elde edilmesi. \\
    \addlinespace
    Hafta 3-5 & Kapsamlı özellik mühendisliği: Bölüm \ref{sec:approach}'ta belirtilen hibrit benzerlik skorları dahil tüm özelliklerin kodlanması. \\
    \addlinespace
    Hafta 6 & Model eğitimi: Oluşturulan zengin özellik tablosu üzerinde farklı makine öğrenmesi modelleri ile denemeler yapılması. \\
    \addlinespace
    Hafta 7 & Özellik analizi: Türetilen özelliklerin model başarısına (pozitif/negatif) etkisinin nicel olarak incelenmesi. \\
    \addlinespace
    Hafta 8 & Proje raporunun (IEEE formatında) hazırlanması ve teslimi. \\
    \bottomrule
  \end{tabular}
\end{table}

\subsection{B Planı}
\label{subsec:planb}

Bu projedeki en zaman alıcı ve teknik olarak zorlu adım, Bölüm \ref{sec:approach}'de ``Hibrit Kullanıcı-Makale Etkileşim Özellikleri" (örn. kosinüs benzerliği) olarak tanımlanan özelliklerin türetilmesidir.

\textbf{B Planı:} Eğer bu gelişmiş hibrit özelliklerin (kullanıcı profili vektörleri vb.) implementasyonu beklenen sürede tamamlanamazsa veya bu özellikler model performansına anlamlı bir katkı sağlamazsa, farklı veri yapıları ile kullanıcı-makale ilişkisi modellenmeye çalışılabilir: İkili çizge yapıları ya da ikili ilişkileri modellemede kullanılan farklı metotlar denenebilir. Bu tip özellik vektörleriyle başarılı sonuçlar alınmaması durumunda text verisi (haber başlık ve gövdeleri) üzerinden özellik vektörü çıkarımı yapılmaya çalışılabilir.


% An example of a floating figure using the graphicx package.
% Note that \label must occur AFTER (or within) \caption.
% For figures, \caption should occur after the \includegraphics.
% Note that IEEEtran v1.7 and later has special internal code that
% is designed to preserve the operation of \label within \caption
% even when the captionsoff option is in effect. However, because
% of issues like this, it may be the safest practice to put all your
% \label just after \caption rather than within \caption{}.
%
% Reminder: the "draftcls" or "draftclsnofoot", not "draft", class
% option should be used if it is desired that the figures are to be
% displayed while in draft mode.
%
%\begin{figure}[!t]
%\centering
%\includegraphics[width=2.5in]{myfigure}
% where an .eps filename suffix will be assumed under latex, 
% and a .pdf suffix will be assumed for pdflatex; or what has been declared
% via \DeclareGraphicsExtensions.
%\caption{Simulation results for the network.}
%\label{fig_sim}
%\end{figure}

% Note that the IEEE typically puts floats only at the top, even when this
% results in a large percentage of a column being occupied by floats.


% An example of a double column floating figure using two subfigures.
% (The subfig.sty package must be loaded for this to work.)
% The subfigure \label commands are set within each subfloat command,
% and the \label for the overall figure must come after \caption.
% \hfil is used as a separator to get equal spacing.
% Watch out that the combined width of all the subfigures on a 
% line do not exceed the text width or a line break will occur.
%
%\begin{figure*}[!t]
%\centering
%\subfloat[Case I]{\includegraphics[width=2.5in]{box}%
%\label{fig_first_case}}
%\hfil
%\subfloat[Case II]{\includegraphics[width=2.5in]{box}%
%\label{fig_second_case}}
%\caption{Simulation results for the network.}
%\label{fig_sim}
%\end{figure*}
%
% Note that often IEEE papers with subfigures do not employ subfigure
% captions (using the optional argument to \subfloat[]), but instead will
% reference/describe all of them (a), (b), etc., within the main caption.
% Be aware that for subfig.sty to generate the (a), (b), etc., subfigure
% labels, the optional argument to \subfloat must be present. If a
% subcaption is not desired, just leave its contents blank,
% e.g., \subfloat[].


% An example of a floating table. Note that, for IEEE style tables, the
% \caption command should come BEFORE the table and, given that table
% captions serve much like titles, are usually capitalized except for words
% such as a, an, and, as, at, but, by, for, in, nor, of, on, or, the, to
% and up, which are usually not capitalized unless they are the first or
% last word of the caption. Table text will default to \footnotesize as
% the IEEE normally uses this smaller font for tables.
% The \label must come after \caption as always.
%
%\begin{table}[!t]
%% increase table row spacing, adjust to taste
%\renewcommand{\arraystretch}{1.3}
% if using array.sty, it might be a good idea to tweak the value of
% \extrarowheight as needed to properly center the text within the cells
%\caption{An Example of a Table}
%\label{table_example}
%\centering
%% Some packages, such as MDW tools, offer better commands for making tables
%% than the plain LaTeX2e tabular which is used here.
%\begin{tabular}{|c||c|}
%\hline
%One & Two\\
%\hline
%Three & Four\\
%\hline
%\end{tabular}
%\end{table}


% Note that the IEEE does not put floats in the very first column
% - or typically anywhere on the first page for that matter. Also,
% in-text middle ("here") positioning is typically not used, but it
% is allowed and encouraged for Computer Society conferences (but
% not Computer Society journals). Most IEEE journals/conferences use
% top floats exclusively. 
% Note that, LaTeX2e, unlike IEEE journals/conferences, places
% footnotes above bottom floats. This can be corrected via the
% \fnbelowfloat command of the stfloats package.




% conference papers do not normally have an appendix


% use section* for acknowledgment



% trigger a \newpage just before the given reference
% number - used to balance the columns on the last page
% adjust value as needed - may need to be readjusted if
% the document is modified later
%\IEEEtriggeratref{8}
% The "triggered" command can be changed if desired:
%\IEEEtriggercmd{\enlargethispage{-5in}}

% references section

% can use a bibliography generated by BibTeX as a .bbl file
% BibTeX documentation can be easily obtained at:
% http://mirror.ctan.org/biblio/bibtex/contrib/doc/
% The IEEEtran BibTeX style support page is at:
% http://www.michaelshell.org/tex/ieeetran/bibtex/
%\bibliographystyle{IEEEtran}
% argument is your BibTeX string definitions and bibliography database(s)
%\bibliography{IEEEabrv,../bib/paper}
%
% <OR> manually copy in the resultant .bbl file
% set second argument of \begin to the number of references
% (used to reserve space for the reference number labels box)
\begin{thebibliography}{1}


\bibitem{kruse2024recsys}
J. Kruse, K. Lindskow, A. Uppal, M. R. Andersen, J. Frellsen, M. Polignano, C. Pomo, ve A. Srivastava, ``RecSys Challenge 2024: Balancing Accuracy and Editorial Values in News Recommendations," \emph{arXiv preprint arXiv:2409.20483}, 2024.

\bibitem{wu2019npa}
C. Wu, F. Wu, S. Yang, Y. Huang, ve X. Xie, ``NPA: Neural news recommendation with personalized attention," in \emph{Proc. ACM SIGKDD Int. Conf. on Knowledge Discovery \& Data Mining (KDD)}, 2019, ss. 2576--2584.

\bibitem{an2019lstur}
M. An, F. Wu, C. Wu, T. Zhang, Y. Xie, ve X. Xie, ``Neural news recommendation with long- and short-term user representations," in \emph{Proc. Annual Meeting of the Association for Computational Linguistics (ACL)}, 2019, ss. 116--125.

\bibitem{hidasi2015gru4rec}
B. Hidasi, A. Karatzoglou, L. Baltrunas, ve D. Tikk, ``Session-based recommendations with recurrent neural networks," in \emph{Proc. Int. Conf. on Learning Representations (ICLR)}, 2016.

\bibitem{kang2018sasrec}
W. C. Kang ve J. McAuley, ``Self-attentive sequential recommendation," in \emph{Proc. IEEE Int. Conf. on Data Mining (ICDM)}, 2018, ss. 197--206.



\end{thebibliography}




% that's all folks
\end{document}


